\documentclass{beamer}
\mode<presentation>
\usepackage{amsmath}
\usepackage{amssymb}
\usepackage{bm}


%\usepackage{advdate}
\usepackage{adjustbox}
\usepackage{subcaption}
%\usepackage{enumitem}
\usepackage{enumerate}
\usepackage{multicol}
\usepackage{mathtools}
\usepackage{listings}
\usepackage{url}
\def\UrlBreaks{\do\/\do-}
\usetheme{Boadilla}
\usecolortheme{lily}
\setbeamertemplate{footline}
{
  \leavevmode%
  \hbox{%
  \begin{beamercolorbox}[wd=\paperwidth,ht=2.25ex,dp=1ex,right]{author in head/foot}%
    \insertframenumber{} / \inserttotalframenumber\hspace*{2ex} 
  \end{beamercolorbox}}%
  \vskip0pt%
}
\setbeamertemplate{navigation symbols}{}

\providecommand{\nCr}[2]{\,^{#1}C_{#2}} % nCr
\providecommand{\nPr}[2]{\,^{#1}P_{#2}} % nPr
\providecommand{\mbf}{\mathbf}
\providecommand{\pr}[1]{\ensuremath{\Pr\left(#1\right)}}
\providecommand{\qfunc}[1]{\ensuremath{Q\left(#1\right)}}
\providecommand{\sbrak}[1]{\ensuremath{{}\left[#1\right]}}
\providecommand{\lsbrak}[1]{\ensuremath{{}\left[#1\right.}}
\providecommand{\rsbrak}[1]{\ensuremath{{}\left.#1\right]}}
\providecommand{\brak}[1]{\ensuremath{\left(#1\right)}}
\providecommand{\lbrak}[1]{\ensuremath{\left(#1\right.}}
\providecommand{\rbrak}[1]{\ensuremath{\left.#1\right)}}
\providecommand{\cbrak}[1]{\ensuremath{\left\{#1\right\}}}
\providecommand{\lcbrak}[1]{\ensuremath{\left\{#1\right.}}
\providecommand{\rcbrak}[1]{\ensuremath{\left.#1\right\}}}
\providecommand{\rank}{\text{rank}}
\theoremstyle{remark}
\newtheorem{rem}{Remark}
\newcommand{\sgn}{\mathop{\mathrm{sgn}}}
\providecommand{\abs}[1]{\left\vert#1\right\vert}
\providecommand{\res}[1]{\Res\displaylimits_{#1}} 
\providecommand{\norm}[1]{\lVert#1\rVert}
\providecommand{\mtx}[1]{\mathbf{#1}}
\providecommand{\mean}[1]{E\left[ #1 \right]}
\providecommand{\fourier}{\overset{\mathcal{F}}{ \rightleftharpoons}}
%\providecommand{\hilbert}{\overset{\mathcal{H}}{ \rightleftharpoons}}
\providecommand{\system}{\overset{\mathcal{H}}{ \longleftrightarrow}}
	%\newcommand{\solution}[2]{\vec{Solution:}{#1}}
%\newcommand{\solution}{\noindent \vec{Solution: }}
\providecommand{\dec}[2]{\ensuremath{\overset{#1}{\underset{#2}{\gtrless}}}}
\newcommand{\myvec}[1]{\ensuremath{\begin{pmatrix}#1\end{pmatrix}}}
\newenvironment{amatrix}[1]{%
  \left(\begin{array}{@{}*{#1}{c}|c@{}}
}{%
  \end{array}\right)
}
\let\vec\mathbf

\lstset{
%language=C,
frame=single, 
breaklines=true,
columns=fullflexible
}

%\numberwithin{equation}{section}

\title{Matgeo-2.9.5}
\author{S.Harsha Vardhan Reddy-ee25btech11054}
\date{\today} 
\begin{document}

\begin{frame}
\titlepage
\end{frame}

\section*{Outline}

\begin{frame}
\frametitle{Question}

\par
$|\vec{a}| = 3, |\vec{b}| = 2\sqrt{3}, \text{ and } \vec{a}\cdot\vec{b} = 6$, then find the value of $|\vec{a}\times\vec{b}|$. \par
\end{frame}
\begin{frame}

\begin{solution}
\par

Given:
\begin{align}
 \norm{\vec{a}} &= 3 \\
 \norm{\vec{b}} &= 2\sqrt{3} \\
 \vec{a}^\top\vec{b} &= 6
\end{align}
We know:
\begin{align}
 \norm{\vec{a}\times\vec{b}} &= \norm{\vec{a}}\norm{\vec{b}}\sin\theta \\
 \cos\theta &= \frac{\vec{a}^\top\vec{b}}{\norm{\vec{a}}\norm{\vec{b}}}
\end{align}
\end{solution}
\end{frame}
\begin{frame}
\begin{solution}
    Thus
\begin{align}
 \brak{\vec{a}^\top\vec{b}}^2 + \brak{\norm{\vec{a}\times\vec{b}}}^2 = \norm{\vec{a}}^2\norm{\vec{b}}^2
\end{align}
Substituting values
\begin{align}
\norm{\vec{a}\times\vec{b}} &= \sqrt{(3^2)\times (2\sqrt{3})^2 - 6^2} \\
\norm{\vec{a}\times\vec{b}} &= \sqrt{108-36} \\
\norm{\vec{a}\times\vec{b}} &=  \sqrt{72} \\
\norm{\vec{a}\times\vec{b}} &=  6\sqrt{2} \\
\end{align}
\end{solution}
\end{frame}
\begin{frame}{Plot}
\begin{figure}
    \centering
    \includegraphics[width=0.8\columnwidth]{figs/plot.png}
    \caption{}
    \label{fig:placeholder}
\end{figure}
\end{frame}



\end{document}
